\documentclass[10pt]{article}
\usepackage[T1]{fontenc}
\usepackage[utf8]{inputenc}
\usepackage[french]{babel}
\usepackage{amsmath,amssymb,mathrsfs,mathpazo}
\usepackage{geometry}
\geometry{a4paper,margin=2cm}
\usepackage{graphicx,array,tabularx,tabularray,colortbl,multirow,multicol,booktabs,enumitem,pifont,ulem,fancybox,diagbox,cancel}
\usepackage{pstricks,pstricks-add,pst-plot,pst-tree,pst-node,pst-text,pst-func,pst-math,pst-bspline,pst-3dplot,pst-eucl,pst-solides3d}
\let\pscalculate\undefined
\usepackage{pst-bezier}
\usepackage{tikz,pgfplots,tkz-tab}
\usetikzlibrary{arrows,arrows.meta,patterns,calc,positioning,trees,decorations.markings,shapes.geometric,shapes.misc,shapes.arrows,intersections,decorations.pathmorphing,decorations.pathreplacing,backgrounds,shadows,babel}
\pgfplotsset{compat=1.18}
\usepgfplotslibrary{fillbetween}
\usepackage[np]{numprint}
\usepackage{listings,makecell,pas-tableur,twoopt}
\let\barre\undefined
\usepackage[pstricks]{tablvar}
\definecolor{darkgreen}{RGB}{0,100,0}
\providecommand{\red}{\color{red}}
\providecommand{\blue}{\color{blue}}
\providecommand{\green}{\color{green}}
\providecommand{\black}{\color{black}}
\providecommand{\white}{\color{white}}
\providecommand{\gray}{\color{gray}}
\providecommand{\cyan}{\color{cyan}}
\providecommand{\magenta}{\color{magenta}}
\providecommand{\yellow}{\color{yellow}}
\providecommand{\violet}{\color{violet}}
\providecommand{\orange}{\color{orange}}
\providecommand{\euro}{\texteuro}
\providecommand{\EUR}{\texteuro}
\providecommand{\ieme}{\textsuperscript{e}}
\providecommand{\degres}{\ensuremath{^\circ}}
\providecommand{\degree}{\ensuremath{^\circ}}
\providecommand{\vv}[1]{\overrightarrow{#1}}
\providecommand{\Coord}[1]{#1}
\providecommand{\decosix}{}
\providecommand{\starredbullet}{\ensuremath{\bullet}}
\pagestyle{empty}
\setlength{\parindent}{0pt}
\usepackage{adjustbox}\DeclareUnicodeCharacter{2236}{:}\DeclareUnicodeCharacter{2190}{\ensuremath{\leftarrow}}
\providecommand{\R}{\mathbb{R}}
\providecommand{\N}{\mathbb{N}}
\providecommand{\Z}{\mathbb{Z}}
\providecommand{\Q}{\mathbb{Q}}
\providecommand{\D}{\mathbb{D}}
\providecommand{\C}{\mathbb{C}}
\providecommand{\e}{{\mathrm{e}}}
\providecommand{\vect}[1]{\overrightarrow{#1}}
\providecommand{\vectt}[1]{\overrightarrow{\mathrm{#1}}}
\providecommand{\vvt}[1]{\overrightarrow{\mathrm{#1}}}
\providecommand{\np}[2][]{{#2~#1}}
\providecommand{\vv}[1]{{\overrightarrow{#1}}}
\renewcommand{\d}{{\mathrm{d}}}
\providecommand{\degres}{{^\circ}}
\providecommand{\degre}{{^\circ}}
\providecommand{\blue}{\color{blue}}
\providecommand{\red}{\color{red}}
\providecommand{\enskip}{\quad}
\providecommand{\llbracket}{\lbrack\!\lbrack}
\providecommand{\rrbracket}{\rbrack\!\rbrack}
\providecommand{\euro}{{\text{€}}}
\providecommand{\no}{{\text{n°}}}
\providecommand{\interval}[2]{{[#1;#2]}}
\providecommand{\intervFO}[2]{{[#1;#2[}}
\providecommand{\dotfill}{\ldots}

\providecommand{\ds}{\displaystyle}
\providecommand{\textding}[1]{#1}
\providecommand{\Oij}{\(\mathrm{O};\vec{\imath},\vec{\jmath}\)}
\providecommand{\Oijk}{\(\mathrm{O};\vec{\imath},\vec{\jmath},\vec{k}\)}
\providecommand{\Ouv}{\(\mathrm{O};\vec{u},\vec{v}\)}
\providecommand{\intervFF}[2]{\ensuremath{[#1;#2]}}
\renewcommand{\euro}{\ifmmode\text{\texteuro}\else\texteuro\fi}
\DeclareUnicodeCharacter{00B0}{\ifmmode^{\circ}\else\textdegree\fi}
\DeclareUnicodeCharacter{20AC}{\ifmmode\text{\texteuro}\else\texteuro\fi}
% Énoncés : Éducation nationale. Transcription et corrigés : APMEP, Denis Vergès et contributeurs. Découpage et classement : CoopMaths, Sébastien Lozano, Éric Elter et contributeurs. Adaptation de la présentation : Maths Explorer.
% CC BY-SA : https://coopmaths.fr/www/about/
% https://coopmaths.fr/alea/static/e3c/2020/tex/e3c_2020_00_sujet40_1_cor.tex
\begin{document}
\noindent\textit{APMEP / CoopMaths. CC BY-SA. Présentation adaptée par Maths Explorer.}\par\medskip
\subsection*{Question 1}

Pour tout réel \(x\), on a :
\[
\sin(7\pi - x) = \sin(7\pi - x - 6\pi) = \sin(\pi - x) = \sin(x).
\]
Réponse : \(\textbf{a.}\)

\subsection*{Question 2}

Avec \( A(-3\,;\,5) \), \( B(2\,;\,-2) \) et \( C(1\,;\,7) \), on a :
\[
\overrightarrow{AB} \begin{pmatrix} 5 \\ -7 \end{pmatrix} \quad \text{et} \quad \overrightarrow{AC} \begin{pmatrix} 4 \\ 2 \end{pmatrix}.
\]
D'où :
\[
\overrightarrow{AB} \cdot \overrightarrow{AC} = 5 \times 4 + (-7) \times 2 = 20 - 14 = 6.
\]
Réponse : \(\textbf{b.}\)

\subsection*{Question 3}

On considère la fonction \(f(x) = \dfrac{3x + 4}{x^2 + 1}\).

Sa dérivée est :
\[
f'(x) = \dfrac{3(x^2 + 1) - 2x(3x + 4)}{(x^2 + 1)^2} = \dfrac{3x^2 + 3 - 6x^2 - 8x}{(x^2 + 1)^2} = \dfrac{-3x^2 - 8x + 3}{(x^2 + 1)^2}.
\]
Réponse : \(\textbf{c.}\)

\subsection*{Question 4}

L'ensemble des points \( M(x\,;\,y) \) tels que \( x^2 + y^2 - 10x + 6y + 30 = 0 \) est :

\begin{align*}
&x^2 + y^2 - 10x + 6y + 30 = 0 \\
\iff &(x - 5)^2 - 25 + (y + 3)^2 - 9 + 30 = 0 \\
\iff &(x - 5)^2 + (y + 3)^2 - 4 = 0 \\
\iff &(x - 5)^2 + (y + 3)^2 = 2^2,
\end{align*}
on reconnaît une équation du cercle de centre \( (5\,;\,-3) \) et de rayon 2.

Réponse : \(\textbf{c.}\)

\subsection*{Question 5}

La somme \( 1 + 5 + 5^2 + \dots + 5^{30} \) est égale à :

\[
S= 1 + 5 + 5^2 + 5^3 + \dots + 5^{30},
\]
donc :
\[
5S = 5 + 5^2 + 5^3 + \dots + 5^{30} + 5^{31},
\]
et par différence :
\[
4S = 5^{31} - 1 \iff S = \dfrac{5^{31} - 1}{4}.
\]

Réponse : \(\textbf{d.}\)

\end{document}
